\documentclass[11pt]{article}
\usepackage[utf8]{inputenc}
\usepackage[T1]{fontenc}
\usepackage{amsmath,amssymb}
\usepackage[margin=1.1in]{geometry}
\usepackage{booktabs}
\usepackage{authblk}
\usepackage[colorlinks=true,linkcolor=blue,citecolor=blue,urlcolor=blue]{hyperref}

\title{The Illusion of Precision in Quantum Mechanics}

\author[1]{Claes Johnson}
\affil[1]{KTH Royal Institute of Technology, SE-100 44 Stockholm, Sweden \\
\small\texttt{claesjohnson@gmail.com}}

\date{\today}

\begin{document}
\maketitle

\begin{abstract}
A fitted model is exact in the aspect it was fitted to. That exactness is automatic and carries no
information, because it was arranged; what a model does \emph{outside} its fit is the only thing
that can fail, and therefore the only thing that is evidence. Textbook quantum mechanics begins by
reporting the hydrogen ground state as $-0.500000$ and treating this as exact, then concludes that
the Schr\"odinger equation is exact for any number of electrons and merely uncomputable. Both steps
deserve examination. The figure $-0.500000$ mixes a prediction with a convention: the pattern
$1/n^2$ is predicted, while the coefficient $\tfrac12$ defines atomic units, and the bridge from
atomic energy to measurable frequency is Planck's constant, itself fixed by the spectra it is then
used to reproduce. The bound-state spectrum is thus the fitted aspect, and exactness there is free.
The first genuinely unfitted predictions fail: applied correctly as a two-body problem, which it
handles exactly, the equation reaches only $1.1\times10^{-5}$ on the ground state, and it makes
$2s$ and $2p$ degenerate against a measured splitting of $1058$~MHz. It is the leading term of an
expansion completed by Dirac and by quantum electrodynamics. The $3N$-dimensional equation for $N$
electrons is then a separate postulate, never solved without approximation beyond three or four
electrons, so its exactness is supported only by extrapolation from a case where exactness is
tautological. The electron's anomalous moment to twelve digits, usually offered in evidence, is a
one-particle vacuum quantity of quantum electrodynamics and says nothing about it. We set against
this a model of non-overlapping electron charge densities in three dimensions --- equally a
postulate, with no claim to exactness --- whose energies, geometries and ionization energies come
out within a few per cent from one electron to the ten thousand of a protein, parameter-free and
computable. The comparison is not accuracy against inaccuracy, but a precision confined to
few-electron systems against an accuracy that can be stated at every size.
\end{abstract}

\section{The premise, checked}\label{sec:premise}

The claim that organises every introduction to the subject is that the Schr\"odinger equation
solves hydrogen exactly. Written for one electron about a fixed point charge,
\begin{equation}
-\tfrac12\nabla^2\psi - \frac{1}{r}\,\psi = E\psi,
\label{eq:schrH}
\end{equation}
it has the ground state $E = -\tfrac12$ in atomic units, which is $13.605693$~eV, and a spectrum
$E_n = -1/2n^2$ reproducing the Balmer series. The agreement was decisive in 1926 and remains the
reason the equation is believed.

It is not, however, exact, and the inexactness must be identified carefully, because one
candidate explanation is not a defect of the equation at all. The nucleus moves, and
\eqref{eq:schrH} holds it fixed --- but that is a shortcoming of how the equation has been
\emph{applied}, not of the equation. Written for two bodies it separates exactly into
centre-of-mass and relative motion, and the relative part returns the hydrogenic spectrum with the
reduced mass $\mu = m_em_p/(m_e+m_p)$ in place of $m_e$, giving $13.598287$~eV. That is the
equation's own answer, and it is obtained exactly.

The residual against the measured $13.598435$~eV~\cite{NIST} is
\[
0.000147~\text{eV}, \qquad 1.1\times10^{-5}\ \text{relative},
\]
and this one is a genuine deficiency: it is the relativistic and radiative physics the equation
does not contain, supplied by the Dirac equation and by quantum electrodynamics.

The spectrum fails more sharply than the energy, and qualitatively rather than numerically.
Equation~\eqref{eq:schrH} makes $2s$ and $2p$ exactly degenerate. They are not: the splitting is $1057.845$~MHz, the Lamb
shift~\cite{LambRetherford}, whose discovery in 1947 initiated renormalised quantum
electrodynamics. Fine structure enters at order $\alpha^2 \approx 5\times10^{-5}$ of the binding
energy and the Lamb shift at order $\alpha^3 \approx 4\times10^{-7}$; the Schr\"odinger equation
contains neither.

So the equation that is said to be exact for hydrogen is the leading term of an expansion. Its
accuracy for the ground-state energy is five significant figures, not infinitely many, and it
predicts a degeneracy that does not exist. This is well known in atomic physics and is not a
criticism of the equation. It is a correction to what the equation is taken to establish.

\section{Where the form came from}\label{sec:form}

Equation~\eqref{eq:schrH} has two ingredients. The potential term is the Coulomb interaction,
taken from electrostatics and therefore empirical. The kinetic term is the quadratic form
\begin{equation}
T[\psi] \;=\; \tfrac12\int |\nabla\psi|^2\,dx ,
\label{eq:kinetic}
\end{equation}
and this is a choice. It was not derived from anything prior; Schr\"odinger arrived at it by
analogy with the Hamilton--Jacobi theory and de Broglie's relation, and it was adopted because the
resulting eigenvalues matched the hydrogen spectrum. That is model selection by agreement with
measurement, which is the ordinary way physics proceeds and is not a defect.

It is worth separating what is predicted here from what is arranged, because the famous
$-0.500000$ consists of both.

The physical content is the \emph{pattern} $E_n \propto 1/n^2$. That is a prediction, it is not
trivial, and it is what made the equation believable in 1926.

The \emph{number} is not. The coefficient $\tfrac12$ in \eqref{eq:kinetic}, and with it a ground
state at exactly $-\tfrac12$, is a statement about units: atomic units are defined so that the
hydrogen ground state comes out $-\tfrac12$, which makes that figure a convention rather than an
outcome. Converting it into anything measurable requires a bridge from atomic energy to the energy
of light, and that bridge is Planck's constant, through $\Delta E = h\nu$. But $h$ ---
equivalently $R_\infty$, since the Hartree is $2R_\infty hc$ --- is itself fixed by spectroscopy,
and the hydrogen spectrum is among the measurements that fix it. The scale is read off the spectrum
and then put back into it.

Taken together: the choice of \eqref{eq:kinetic} among the quadratic forms one might write, and
the constant connecting its eigenvalues to observed frequencies, were both settled by agreement
with measurement. In the narrow sense the equation has no adjustable parameter --- given $\hbar$,
$m_e$ and $e$ the ground state follows. In the sense that matters for the word \emph{exact}, it is
a model of simple form whose form and whose scale were selected to fit.

And this is the point on which the whole question turns, so it is worth stating as a principle:

\begin{quote}
\emph{A fitted model is exact in the aspect it was fitted to.}
\end{quote}

That exactness is automatic. It carries no information, because it was arranged. The informative
question is never how well a model reproduces what selected it, but what it does \emph{outside}
the fit --- and the answer there is a prediction, which can fail.

For hydrogen the fitted aspect is the bound-state spectrum: the $1/n^2$ pattern and, through $h$,
its scale. Exactness there is free. The first genuinely unfitted predictions are the ones named in
Section~\ref{sec:premise} --- the fine structure and the $2s$--$2p$ splitting --- and the equation
gets both wrong, the second qualitatively, requiring Dirac and then quantum electrodynamics to
repair. That is the ordinary and respectable fate of a good model tested beyond its fit. It is not
the record of an exact equation.

\section{The step that is not justified}\label{sec:step}

For $N$ electrons the equation is written in $3N$ dimensions,
\begin{equation}
\Big[-\tfrac12\sum_{i}\nabla_i^2 - \sum_i\frac{Z}{r_i} + \sum_{i<j}\frac{1}{r_{ij}}\Big]
\Psi(x_1,\dots,x_N) = E\,\Psi ,
\label{eq:schrN}
\end{equation}
with $\Psi$ antisymmetric. This is a separate postulate. It is not derived from
\eqref{eq:schrH}; it is \eqref{eq:schrH} generalised by writing the same two ingredients once per
electron and adding pairwise repulsion, and it carries two further components ---
antisymmetry and spin --- that one-electron hydrogen never required.

The generalisation is natural and may well be right. What is not available is evidence that it is
\emph{exact}, and the reason is that nobody has solved it. Direct solution of \eqref{eq:schrN}
costs exponentially in $N$ and is out of reach beyond three or four electrons. Every number the
standard formalism produces for a real molecule comes from an approximation --- a basis set, a
functional, an active space, a truncated expansion --- and the accuracy one can state depends on
which approximation was chosen rather than on the equation.

So the claim ``exact in principle, merely uncomputable'' has an unusual structure. It cannot be
tested, because testing it would require the computation that is said to be impossible. And the
inductive support offered for it is the hydrogen case --- where, by Section~\ref{sec:form}, the
exactness is in the fitted aspect and therefore empty, and where, by Section~\ref{sec:premise},
the first unfitted predictions fail and the spectrum is qualitatively wrong about a degeneracy. An extrapolation from a case that does not hold to cases
that cannot be checked is not a strong position, however widely it is held.

\section{The showcase, and what it shows}\label{sec:ae}

The figure usually produced at this point is the electron's anomalous magnetic moment. Theory and
experiment agree on $a_e = (g-2)/2$ to roughly twelve significant
figures~\cite{Fan2023,Aoyama2019}, and this is regularly described as making quantum mechanics the
most precisely tested theory in science. The description is fair; the inference drawn from it is
not.

The quantity is one-particle. It concerns a single electron's interaction with the quantised
electromagnetic vacuum, computed as a perturbative series in $\alpha$ within quantum
electrodynamics. It involves no many-electron wavefunction, no antisymmetry among several
electrons, no $3N$-dimensional equation. It is evidence --- overwhelming evidence --- that QED
describes the vacuum coupling of one charge. It is not evidence about \eqref{eq:schrN}, and
offering it in support of the many-electron equation confuses two different theories with one
Lagrangian between them.

Two further points are worth making precisely. The comparison requires $\alpha$, and the most
precise determinations of $\alpha$ now come either from $a_e$ itself or from atom-recoil
measurements whose results have at times disagreed at the level of the
comparison~\cite{Morel2020,Parker2018}; so part of the agreement is internal rather than
independent. And the precision is a property of the observable, not of the theory's reach:
twelve digits on a vacuum quantity coexists with no ability to state an error bar on a
hundred-electron molecule.

\section{The accuracy that collapses}\label{sec:curve}

Set out the standard formalism's precision as a function of electron count, and the shape is the
thing to notice.

\begin{center}
\begin{tabular}{llll}
\toprule
system & electrons & standard QM & continuum model \\
\midrule
H          & 1    & analytic, $4$ figures vs measurement & exact ($-0.5$ by construction) \\
He         & 2    & $\sim\!20$ digits~\cite{Nakashima2007} & $2$--$4\%$ \\
H$_2$      & 2    & $\sim\!12$ digits~\cite{KolosWolniewicz} & $8.6\%$ total, $2.3\%$ bond length \\
atoms to Ar& 2--18& $\sim\!1$ kcal/mol, heavy machinery & $1.6$--$4.5\%$ \\
protein    & $\sim\!10^4$ & \textbf{no number} & few per cent, $153$ residues \\
\bottomrule
\end{tabular}
\end{center}

At two electrons the gap is some fifteen orders of magnitude in the standard formalism's favour. At
ten thousand there is no comparison to make, because one column is empty. The twenty digits for
helium are obtained by explicitly correlated expansions in interelectronic coordinates --- methods
that do not scale and are not truncations of any general procedure. They are evidence about the
two-electron problem, which is specially tractable, and not about the theory's reach.

\section{A different generalisation, with no claim to exactness}\label{sec:realqm}

The alternative examined here keeps the two ingredients of \eqref{eq:schrH} --- Coulomb
interaction and the kinetic form \eqref{eq:kinetic} --- and generalises differently. Instead of one
antisymmetric function of $3N$ variables, $N$ charge densities occupy
\emph{non-overlapping domains} in ordinary three-dimensional space, meeting with continuity across
free boundaries whose position the energy determines~\cite{Johnson-RealQM}. Exclusion is realised
geometrically rather than by antisymmetry; there is no spin; $Z$ is the only quantity supplied.

Nothing argues that this is exact, and we do not claim it. It is a postulate, chosen for the same
kind of reason \eqref{eq:kinetic} was chosen: it is the simplest generalisation of the two
ingredients that remains three-dimensional, and it can be computed. Its predictions are therefore
to be checked rather than deduced, and the checking has been done over a wide range:

\begin{itemize}\itemsep2pt
\item total atomic energies within $1$--$2\%$ on average across the first three periods, worst case
  $4.5\%$ at silicon, every atom solved the same way with nothing fitted~\cite{Johnson-atoms};
\item helium at $-2.845$ against $-2.904$, $2.0\%$, on the spherically homogenised solver;
\item H$_2$ at $8.6\%$ on the total energy with the bond length at $1.368$ against $1.401$,
  $2.3\%$, from commensurate separations on a converged mesh;
\item first ionization energies of the alkalis and second of the alkaline earths within $1.4\%$ and
  $4\%$ respectively, from one constant per column, against the measured radius of the ion that
  remains~\cite{Johnson-ionization};
\item the row lengths $2,8,8,18,18,32,32$ and closure at all seven noble gases, from the pair and
  the octet with no $2n^2$ entering;
\item hydrogen-bond geometry and secondary structure in polypeptides of up to $153$ residues,
  carried with their electrons on a laptop.
\end{itemize}

The failures should be stated with the same directness. The H$_2$ \emph{binding} energy comes out
$58\%$ high --- which is the $8.6\%$ error on the total divided by the bond energy being a seventh
of it, not an independent defect, but $58\%$ nonetheless. Ionization energies require a radius
supplied from measured ionic radii, so that result is a correlation between two measured
quantities and not a prediction from $Z$. Proton affinities come out with the wrong sign.
Dispersion-bound systems are outside range. The model is not accurate in the sense the standard
formalism is accurate for helium, and it will not become so.

\section{What the comparison is about}\label{sec:conclusion}

The two models are both postulates. Both generalise Coulomb interaction and a chosen kinetic form
to many electrons, in different ways, and neither generalisation is derived. They differ in what
can be done with them.

One is exact in a limit nobody can compute, is wrong at $5\times10^{-4}$ in the one case where
exactness is claimed, achieves extraordinary precision for one and two electrons by methods that do
not scale, and offers no error bar at all past a few dozen. The other is exact nowhere, agrees with
measurement to a few per cent, and does so at every size from one electron to ten thousand, in
three dimensions, at linear cost, with nothing fitted.

The first is more accurate wherever it can be evaluated. The second is \emph{answerable} wherever
matter exists. A theory whose error is knowable and roughly constant can be held to account at
every scale; a theory whose error is unknowable past four electrons can be held to account only
there. Which property matters more depends on what one wants a model of matter for --- and that is
a question about scope, which is a legitimate scientific question, rather than a question about
precision, which the twelve digits have already settled in a domain where nobody disputes them.

\begin{thebibliography}{99}

\bibitem{NIST} A.~Kramida, Yu.~Ralchenko, J.~Reader and NIST ASD Team, \emph{NIST Atomic Spectra
Database}, National Institute of Standards and Technology, Gaithersburg MD.

\bibitem{LambRetherford} W.~E.~Lamb and R.~C.~Retherford, \emph{Fine structure of the hydrogen
atom by a microwave method}, Phys.\ Rev.\ \textbf{72} (1947) 241--243.

\bibitem{Fan2023} X.~Fan, T.~G.~Myers, B.~A.~D.~Sukra and G.~Gabrielse, \emph{Measurement of the
electron magnetic moment}, Phys.\ Rev.\ Lett.\ \textbf{130} (2023) 071801.

\bibitem{Aoyama2019} T.~Aoyama, T.~Kinoshita and M.~Nio, \emph{Theory of the anomalous magnetic
moment of the electron}, Atoms \textbf{7} (2019) 28.

\bibitem{Morel2020} L.~Morel, Z.~Yao, P.~Clad\'e and S.~Guellati-Kh\'elifa, \emph{Determination of
the fine-structure constant with an accuracy of 81 parts per trillion}, Nature \textbf{588} (2020)
61--65.

\bibitem{Parker2018} R.~H.~Parker, C.~Yu, W.~Zhong, B.~Estey and H.~M\"uller, \emph{Measurement of
the fine-structure constant as a test of the Standard Model}, Science \textbf{360} (2018) 191--195.

\bibitem{Nakashima2007} H.~Nakashima and H.~Nakatsuji, \emph{Solving the Schr\"odinger equation
for helium atom and its isoelectronic ions with the free iterative complement interaction method},
J.\ Chem.\ Phys.\ \textbf{127} (2007) 224104.

\bibitem{KolosWolniewicz} W.~Ko{\l}os and L.~Wolniewicz, \emph{Potential-energy curves for the
$X^1\Sigma_g^+$, $b^3\Sigma_u^+$ and $C^1\Pi_u$ states of the hydrogen molecule}, J.\ Chem.\ Phys.\
\textbf{43} (1965) 2429--2441.

\bibitem{Johnson-RealQM} C.~Johnson, \emph{Real Quantum Mechanics}, submitted.

\bibitem{Johnson-atoms} C.~Johnson, \emph{A 3D Multiphase Continuum Model for Atoms and the
Periodic Table}, submitted.

\bibitem{Johnson-ionization} C.~Johnson, \emph{Ionization Energy Without Differencing: Two Columns
from the Radius of the Ion That Remains}, draft.

\end{thebibliography}
\end{document}
